\section*{Acknowledgments}\label{sec:ack}

This work has benefited from the contributions of many individuals across
its development.

The project originated at the Escuela Nacional de Pintura, Escultura y
Grabado ``La Esmeralda'' (ENPEG), where Dra.\ Elena Odgers and Prof.\
Jos\'e Luis S\'anchez Rull provided the intellectual environment and early
guidance that set this research in motion.  At the Instituto Tecnol\'ogico
Aut\'onomo de M\'exico (ITAM), Dr.\ Franz Peter Oberarzbacher Niederwolfsgruber
and Dr.\ Pavel Jim\'enez V\'azquez were instrumental in shaping the
foundational direction of the project.

We owe a particular debt to the mathematician Leonardo Jim\'enez Mart\'inez,
without whom this project would never have developed.

Pablo Tenorio contributed as a collaborator during the earlier stages of the
work, providing independent research support.

We acknowledge Javier Grisales Herrera, whose Golden Prime Symmetry
Theory~\cite{grisales25} provides the cyclotomic classification of primes
essential to the operator construction.

The mathematical landscape of this work was shaped by the foundational
contributions of Yuri Manin, whose $\mathbb{F}_1$-geometric program
inspired the categorical framework, and the pioneering work of Alain Connes
and Matilde Marcolli in noncommutative geometry.  The string-theoretic and
holographic perspectives owe a profound debt to Leonard Susskind, Juan
Maldacena, and Edward Witten.  We are particularly grateful to John Huerta and Urs Schreiber, whose work on the superpoint tower and higher categorical
structures---and whose broader contributions through the $n$Lab
project---provided essential conceptual tools for connecting the
dimensional hierarchy of PCF to string-theoretic frameworks.

We are deeply indebted to Andrew Odlyzko, whose extensive computational verification of Riemann zeros---spanning billions of zeros at extraordinary precision---provided the empirical foundation against which our operator $T^*$ was tested, and whose 1982 correspondence with P\'olya preserved the historical record of the conjecture that motivates this work.

Technical writing and formal verification were assisted by Claude, Gemini, Qwen, MiMo, GLM, Kimi, Composer, Grok together with Lean~4 and Mathlib; we thank the development teams at Anthropic, Alibaba, Anysphere, DeepSeek AI, Google DeepMind, Xiaomi Research, Moonshot AI, Z.ai, MiniMax, xAI, and the Lean Focused Research Organization (FRO). Special recognition is due to Leonardo de Moura and the Lean community for providing the underlying verification infrastructure.

We wish to acknowledge the extensive open-source software ecosystem that facilitated the development of this project. The core computational framework and data visualization components were built using Python, JavaScript, and TypeScript, leveraging numerous libraries for numerical analysis---notably \texttt{mpmath} for high-precision spectral verification---and dynamic chart generation via \texttt{Matplotlib}. Reproducibility is maintained through a deterministic automation pipeline utilizing the \texttt{npm}, \texttt{uv}, and \texttt{PyPI} ecosystems and \texttt{Docker} containerization, integrated with \texttt{Git} (originally developed by Linus Torvalds). This framework is supplemented by \texttt{Google Colab} for cloud-based verification, \texttt{GitHub} (Microsoft) for project hosting, and \texttt{Zenodo} for long-term archiving.

Finally, this work stands on the shoulders of over two millennia of
mathematical thought.  We express our gratitude to every thinker cited in
these pages---from Vitruvius, Euclid, and Cardano through Gauss, Riemann,
and Eisenstein, to Hilbert, P\'olya, G\"odel, Tarski, Grothendieck, and
Weil, and onward to the living contributors whose ideas made this
framework conceivable.  Whatever merit this work may have belongs, in no
small measure, to the cumulative tradition they built.

Science and art have long worked in harmony to revitalize methodologies and chart new intellectual horizons, yielding foundational concepts from modularity to the development of differential geometry and its application to relativity. Inspired by this legacy, and driven by a spirit of inquiry, this work underscores the necessity of multidisciplinary collaboration and the pursuit of complex ideas through diverse methodologies, while leveraging the transformative potential of emerging technologies.
